%-------------------------------%
%  Author: Alessandro Sciarra   %
%    Date: 11 Jul 2025          %
%-------------------------------%

% Partially switch to the green variant of the Z02 theme for this section
% NOTE: This is a low-level patch taking the colours from the
%       theme implementation. I need it only once in this session
%       and therefore I did not look for a more elegant way.
% NOTE: I did not switch completely to not interfere with the color-code
%       of sessions of the rest of the day!
\colorlet{PT}{UnrealFoodPills}
\colorlet{ALERT}{PT}
\colorlet{BGDARK}{GreenBlue}
\colorlet{BGLIGHT}{MicroscopeSlide}
\colorlet{TEXT}{Abyss}

% Branch for git-flow
\newcommand{\mainBranchColor}{PQ}
\newcommand{\supportBranchColor}{Marine}
\newcommand{\branch}[2]{\textcolor{#1}{\texttt{\bfseries #2}}}
\newcommand{\main}{\branch{\mainBranchColor}{main}}
\newcommand{\develop}{\branch{\mainBranchColor}{develop}}
\newcommand{\feature}{\branch{\supportBranchColor}{feature}}
\newcommand{\release}{\branch{\supportBranchColor}{release-*}}
\newcommand{\hotfix}{\branch{\supportBranchColor}{hotfix-*}}

%===============================================================%
\subsection{Git tag}
%~~~~~~~~~~~~~~~~~~~~~~~~~~~~~~~~~~~~~~~~~~~~%
\begin{frame}[fragile]{Tagging the history of a repository}
    \vspace{-8mm}
    \begin{overlayarea}{\textwidth}{0.7\textheight}
        \begin{varblock}{example}[0.96\textwidth]{Lightweight tags}<only@1-2,4>
            They simply are a name for an object (often a commit) and they are usually meant for private or temporary use
        \end{varblock}
        \begin{onlyenv}<1>
            \begin{lstlisting}[style=MyBash]
                $ git tag v1.0
                $ git tag
                |+v1.0+|
                $ git show v1.0
                |+commit aa879f463acd41fc38c7e96090cc1eea279304df
                Author: Alessandro Sciarra <sciarra@itp.uni-frankfurt.de>
                Date:   Wed Aug 16 14:01:08 2025 +0200

                Coffee machine production ready+|

                # Commit differences
            \end{lstlisting}
        \end{onlyenv}
        \begin{varblock}{example}[0.96\textwidth]{Annotated tags}<only@2,4>
            Tag objects contain a creation date, the tagger name and e-mail, a tagging message, and an optional GnuPG signature
        \end{varblock}
        \begin{varblock}{}[0.7\textwidth]{Tags as commits are by default local}<only@4>
            Use \;\bash|git push --tags|\; to push them
        \end{varblock}
        \begin{onlyenv}<3>
            \begin{lstlisting}[style=MyBash, aboveskip=7mm]
                $ git tag -a v1.0 -m\
                >     "Coffee machine software ready for production"
                $ git show v1.0
                |+tag v1.0
                Tagger: Alessandro Sciarra <sciarra@itp.uni-frankfurt.de>
                Date:   Wed Aug 17 09:12:24 2025 +0200

                Coffee machine software ready for production

                commit aa879f463acd41fc38c7e96090cc1eea279304df
                Author: Alessandro Sciarra <sciarra@itp.uni-frankfurt.de>
                Date:   Wed Aug 16 14:01:08 2025 +0200

                Coffee machine production ready+|

                # Commit differences
            \end{lstlisting}
            \begin{varblock}{alert}[0.6\textwidth]{}
                \alert{Prefer annotated tags for releases!}
            \end{varblock}
        \end{onlyenv}
    \end{overlayarea}
\end{frame}
%~~~~~~~~~~~~~~~~~~~~~~~~~~~~~~~~~~~~~~~~~~~~%
\begin{frame}[fragile]{Tagging a code: Let's speak about the \textbf{same} code!}
    \vspace{-3mm}
    \begin{varblock}{alerted}[0.9\textwidth]{A milestone in development}<only@1>
        A tag in the git history is to some extent a commitment, but it is also a statement of invaluable help about the codebase!
    \end{varblock}
    \begin{varblock}{}[0.9\textwidth]{Why should I?}<only@1>
        \begin{itemize}
            \item The codebase is released
            \item The codebase will be released
            \item The codebase is  private but shared with colleagues
            \item The codebase will be (maybe) inherited
            \item \textbf{The software is used to produce data}
        \end{itemize}
    \end{varblock}
\end{frame}
%===============================================================%


%===============================================================%
\subsection{Semantic versioning}
%~~~~~~~~~~~~~~~~~~~~~~~~~~~~~~~~~~~~~~~~~~~~%
\begin{frame}{How should tags be named? \hfill\uncover<2->{\ttfamily\normalsize\alert{\raisebox{1pt}{[Prefix]X[.Y[.Z]]}}}}
    \vspace{-2mm}
    \PrepareURLsymbol[PB]
    \begin{varblock*}{}[\textwidth]{According to the \URL*[PB]{https://semver.org/spec/v2.0.0.html}{Semantic Versioning}\color{PP}, increase}<uncover@2->
        \small
        \makebox[9mm][r]{\PQ{MAJOR}} version \PT{\texttt{X}} when you make incompatible API changes\\
        \makebox[9mm][r]{\PQ{MINOR}} version \PT{\texttt{Y}} when you add functionality in a backwards-compatible manner\\
        \makebox[9mm][r]{\PQ{PATCH}} version \PT{\texttt{Z}} when you make backwards-compatible bug fixes
    \end{varblock*}
    \begin{varblock*}{example}[0.95\textwidth]{Choose your alternative, but give yourself a rule, e.g. increase}<uncover@3->
        \small
        \makebox[9mm][r]{\PQ{MAJOR}} version \PT{\texttt{X}} when you introduce new big features\\
        \makebox[9mm][r]{\PQ{MINOR}} version \PT{\texttt{Y}} when you add minor functionality or do big refactoring\\
        \makebox[9mm][r]{\PQ{PATCH}} version \PT{\texttt{Z}} when you make bug fixes (without large changes)
    \end{varblock*}
    \vspace{-2mm}
    \begin{varblock*}{example}[0.95\textwidth]{}<uncover@4->
        \small
        \makebox[9mm][r]{\PQ{MAJOR}} version \PT{\texttt{X}} when you introduce new features\\
        \makebox[9mm][r]{\PQ{MINOR}} version \PT{\texttt{Y}} when you do big refactoring or fix bugs
    \end{varblock*}
    \begin{center}
        \uncover<5>{\alert{However you do it, use a CHANGELOG file to list user-relevant changes!}}
    \end{center}
\end{frame}
%~~~~~~~~~~~~~~~~~~~~~~~~~~~~~~~~~~~~~~~~~~~~%
\begin{frame}[fragile]{Ideally}
    \vspace{-2mm}
    \begin{lstlisting}[style=MyBash]
        $ git tag -n
        |+  v2.0.0   Coffee machine ready for milk drinks      +||-Major-|
        |+  v1.3.1   Fix milk temperature problem              +||*Patch*|
        |+  v1.3.0   Interface milk system with machine        +||=Minor=|
        |+  v1.2.0   Add foam generator                        +||=Minor=|
        |+  v1.1.0   Add milk container                        +||=Minor=|
        |-  v1.0.0-|   |+Coffee machine production ready+|           |-Major-|
        |+  v0.7.1   Fix pipe failures                         +||*Patch*|
        |+  v0.7.0   Interface all components                  +||=Minor=|
        |+  v0.6.0   Add coffee grounds container              +||=Minor=|
        |+  v0.5.0   Add water tray                            +||=Minor=|
        |+  v0.4.0   Add water tank                            +||=Minor=|
        |+  v0.3.0   Add bean container                        +||=Minor=|
        |+  v0.2.0   Add brew system and core engine           +||=Minor=|
        |+  v0.1.0   Deploy coffee machine skeleton            +||=Minor=|
    \end{lstlisting}
    \begin{itemize}
        \item Clear evolution of the codebase for everybody
        \item \textbf{v1.0} usually refers to the first \textbf{production-ready} version
    \end{itemize}
\end{frame}
%===============================================================%


%===============================================================%
\subsection{Git-flow}
%~~~~~~~~~~~~~~~~~~~~~~~~~~~~~~~~~~~~~~~~~~~~%
\begin{frame}<beamer>{\\ \colorbox{PP}{\makebox[1em]{\textcolor{BGLIGHT}{1\vphantom{$A^{\int}_x$}}}}\;Git in real life}
    \vspace{3mm}
    \tableofcontents[sections={1}, currentsection, hideothersubsections, sectionstyle=hide, subsectionstyle=show/shaded/hide, subsubsectionstyle=show]
\end{frame}
%---------------------------------------------------------------%
\subsubsection{Motivation and background}
%~~~~~~~~~~~~~~~~~~~~~~~~~~~~~~~~~~~~~~~~~~~~%
\begin{frame}{Disclaimer and references}
    \PrepareURLsymbol[PT]
    \vspace{-3mm}
    \begin{varblock}{alert}[\textwidth]{\URL*[PT]{http://nvie.com/posts/a-successful-git-branching-model/}{Vincent Driessen} back in 2010}
        In the following I will give you a summary of his blog page.\\
        All credit goes to the author and no original part is here contained.
    \end{varblock}
    \vspace{-1mm}
    \begin{varblock}{}[\textwidth]{Vincent Driessen as \;\href{https://github.com/nvie/gitflow}{\faGithub\;\textbf{nvie}}}
        The original author also implemented a collection of git extensions which allow to easily use his branching model (active till 2012).
    \end{varblock}
    \vspace{-1mm}
    \PrepareURLsymbol[PS]
    \begin{varblock}{example}[0.9\textwidth]{Alternative implementations}
        \begin{tabular}{r@{\;\;}l@{\;\;}l}
            AVH Edition: &
            \PS{\href{https://github.com/petervanderdoes/gitflow-avh}{\faGithub\;gitflow-avh}} &
            (archived in June 2023)\\
            CJS Edition: &
            \PS{\href{https://github.com/CJ-Systems/gitflow-cjs}{\faGithub\;gitflow-cjs}} &
            (still active, but rarely)\\
        \end{tabular}
    \end{varblock}
    \FrameRemark{The AVH edition has an associated \URL*[PS]{https://danielkummer.github.io/git-flow-cheatsheet/index.html}{cheatsheet}, which also illustrate the branching model.}\\
\end{frame}
%~~~~~~~~~~~~~~~~~~~~~~~~~~~~~~~~~~~~~~~~~~~~%
\begin{frame}[c]{Motivation}
    \vspace{-4mm}
    \begin{center}
        \begin{tikzpicture}[every node/.style={inner sep=0mm}]
            \begin{scope}[scope on=<1->]
                \node (L) {\includegraphics[height=28mm]{Git-bad-commits}};
                \node[below = 2mm of L, font=\tiny] {\URL*{https://www.xkcd.com/1296/}{https://www.xkcd.com/1296/}};
            \end{scope}
            \begin{scope}[scope on=<2->]
                \node[right = 1mm of L] (R) {\includegraphics[height=28mm]{Git-messy-history}};
                \node[below = 2mm of R, font=\tiny]
                {\URL*{https://image.slidesharecdn.com/gitseries-150826152158-lva1-app6892/95/git-series-episode-3-git-flow-and-githubflow-2-638.jpg?cb=1443100187}{https://image.slidesharecdn.com/[...]}};
                \fill[white] (R.north west) rectangle ($(R.north east)!0.4!(R.east)$);
                \node[font=\footnotesize, anchor=north, fill=white, text=PS] at ($(R.north)-(0,1.5mm)$)
                {Sometimes history becomes messy};
            \end{scope}
        \end{tikzpicture}
    \end{center}
    \vspace{-3mm}
    \begin{uncoverenv}<3>
        \setlength{\leftmargini}{7mm}
        \begin{itemize}
            \item \alert{Ordered} and \alert{standardized} way to daily work
            \item Easy way to keep big projects history and development tidied up
            \item Sustainable work in team also \alert{for new members}
            \item In larger projects, easy interaction within and between sub-teams
            \item \PP{A way to go for released software}
        \end{itemize}
    \end{uncoverenv}
\end{frame}
%~~~~~~~~~~~~~~~~~~~~~~~~~~~~~~~~~~~~~~~~~~~~%
\begin{frame}[c]{Background}
    \vspace{-3mm}
    \begin{columns}[c]
        \begin{column}{0.5\textwidth}
            \small
            \setlength{\leftmargini}{7mm}
            \begin{itemize}
                \setlength{\itemsep}{0mm}
                \item Central repository (e.g.\ GitHub)
                \item Developers pull/push from/to origin
                \item Sub-team fetches lead to high work efficiency
                \item It might be useful to work together with two or more developers on a big new feature,
                before pushing the work in progress to origin prematurely.
                \item \alert{Never forget \texttt{origin} is public!}
            \end{itemize}
        \end{column}
        \begin{column}{0.45\textwidth}
            \vspace{3mm}
            \centering
            \includegraphics[width=\textwidth]{Centr-decentr}
            \begin{block}{}
                \footnotesize\itshape
                Technically, this means nothing more than that, for example, Alice has defined a Git remote,
                named \textnormal{\texttt{bob}}, pointing to Bob's repository, and vice versa.
            \end{block}
        \end{column}
    \end{columns}
\end{frame}
%---------------------------------------------------------------%

%---------------------------------------------------------------%
\subsubsection{The main branches}
%---------------------------------------------------------------%
\begin{frame}[c]{The main and develop branches}
    \small
    \begin{columns}[b]
        \begin{column}{0.5\textwidth}
            \vspace{-3mm}
            \begin{varblock}{}[0.95\textwidth]{\texttt{origin/\main}}
                The source code of \texttt{HEAD} always reflects a production-ready state
            \end{varblock}
            \vspace{-1mm}
            \begin{varblock}{}[0.95\textwidth]{\texttt{origin/\develop}}
                The source code of \texttt{HEAD} always reflects a state
                with the latest delivered development changes for the next release
            \end{varblock}
            \begin{varblock}{alerted}[0.95\textwidth]{}
                \alert{Each time when changes are merged back into main, this is a new production release by definition}\\
                {\scriptsize $\to$ hooks for enforced policies!}
            \end{varblock}
        \end{column}
        \begin{column}{0.4\textwidth}
            \includegraphics[height=0.77\textheight]{Main-branches}
        \end{column}
    \end{columns}
    \FrameRemark{\URL*{https://git-scm.com/book/en/v2/Customizing-Git-Git-Hooks}{Git hooks} are authomatised actions that are triggered by a given git command, i.e.\ git will automatically run a script whenever you run a git command.}
\end{frame}
%---------------------------------------------------------------%

%---------------------------------------------------------------%
\subsubsection{Feature branches}
%---------------------------------------------------------------%
\begin{frame}[fragile,c]{The feature branches}
    \begin{columns}[c]
        \begin{column}{0.65\textwidth}
            \vspace*{-1cm}
            \begin{overlayarea}{\textwidth}{0.75\textheight}
                \begin{varblock}{}[0.9\textwidth]{Properties}<only@1>
                    \begin{itemize}
                        \item May branch off from:\\
                        \quad{\small\develop}
                        \item Must merge back into:\\
                        \quad{\small\develop}
                        \item Any name except:\\
                        \quad{\small\main,\;\develop}\\[-1mm]
                        \quad{\small\release,\;\hotfix}
                    \end{itemize}
                \end{varblock}
                \vspace{-1mm}
                \begin{varblock}{alerted}[0.9\textwidth]{}<only@1>
                    \small \alert{Feature branches typically exist in developer repository only, not in \texttt{origin}.}
                \end{varblock}
                \vspace{3mm}
                \begin{varblock}{}[1.05\textwidth]{How it looks like:}<only@2>
                    \setlength{\medskipamount}{0pt}
                    \begin{lstlisting}[style=MyBash]
                        $ git switch -c myfeature develop
                        |+Switched to a new branch "myfeature"+|
                        #------------------------------------
                        # Some development, commits, etc.
                        #------------------------------------
                        $ git switch develop
                        |+Switched to branch 'develop'+|
                        $ git pull origin develop
                        |+Already up-to-date.+|
                        $ git merge --no-ff myfeature
                        |+Updating ea1b82a..05e9557
                        (Summary of changes)+|
                        $ git branch -d myfeature
                        |+Deleted branch myfeature (was 05e9557).+|
                        $ git push origin develop
                        |+[...]+|
                    \end{lstlisting}
                \end{varblock}
            \end{overlayarea}
        \end{column}
        \begin{column}{0.3\textwidth}
            \centerline{\includegraphics[height=0.75\textheight]{Feature}}
        \end{column}
    \end{columns}
\end{frame}
%~~~~~~~~~~~~~~~~~~~~~~~~~~~~~~~~~~~~~~~~~~~~%
\begin{frame}[c]{Why to avoid a fast-forward merge?}
    \begin{tikzpicture}[remember picture, overlay]
        \node[visible on=<1>] at ([yshift=-5mm]current page.center) {\includegraphics[height=0.65\textheight]{No-ff}};
    \end{tikzpicture}
    \begin{onlyenv}<2>
        \vspace{-8mm}
        \begin{center}
            \includegraphics[height=0.35\textheight]{No-ff}
        \end{center}
        \setlength{\leftmargini}{7mm}
        \begin{itemize}
            \setlength{\itemsep}{2mm}
            \item[\yes] No information lost about the historical existence of a feature
            \item[\yes] Easy to see in history which commits have implemented a feature
            \item[\yes] Easy to revert a whole feature (i.e.\ a group of commits)
            \item[\no]  It will create (empty) commit objects\\
            \then the gain is much bigger than the cost!
        \end{itemize}
    \end{onlyenv}
\end{frame}
%---------------------------------------------------------------%

%---------------------------------------------------------------%
\subsubsection{Release branches}
%---------------------------------------------------------------%
\begin{frame}[fragile,c]{The release branches}
    \begin{tikzpicture}[overlay, remember picture]
        \node[anchor=north east] (yeah) at ([shift={(-8mm,-3mm)}]current page.north east) {
\includegraphics[height=20mm]{Meme-success}};
        \path[to] ([shift={(-8mm,-9mm)}]current page.north) edge[out=0, in=180] node[midway, font=\tiny, sloped, above=-3pt]{Make you feel like} (yeah);
    \end{tikzpicture}
    \begin{onlyenv}<2>
        \begin{columns}[c]
            \begin{column}{0.55\textwidth}
                \begin{varblock}{}[\textwidth]{Preperties:}
                    \begin{itemize}
                        \item May branch off from:\\
                        \quad{\small\develop}
                        \item Must merge back into:\\
                        \quad{\small\develop,\;\main}
                        \item Branch naming convention:\\
                        \quad{\small\release}
                    \end{itemize}
                \end{varblock}
            \end{column}
            \begin{column}{0.45\textwidth}
                \vspace{6mm}
                \includegraphics[width=\textwidth]{Release}
            \end{column}
        \end{columns}
    \end{onlyenv}
    \begin{onlyenv}<3>
        \begin{minipage}[c]{0.95\textwidth}
            \vspace{14mm}
            \setlength{\leftmargini}{1mm}
            \begin{itemize}
                \setlength{\itemsep}{3mm}
                \item Release branches support preparation of a new production release
                \item Last-minute checks, minor bug fixes and preparing meta-data (version number,
                etc.) should be done on a release branch
                \item By doing all of this work on a release branch, the \textbf{develop}
                branch is cleared to receive features for the following release
                \item \alert{When a new \PQ{release} branch is created, all features that are
                    targeted for this release must be merged into the \PQ{develop} branch}
            \end{itemize}
        \end{minipage}
    \end{onlyenv}
\end{frame}
%~~~~~~~~~~~~~~~~~~~~~~~~~~~~~~~~~~~~~~~~~~~~%
\begin{frame}[fragile,c]{The release branch: How it works}
    \begin{overlayarea}{\textwidth}{0.9\textheight}
        \begin{onlyenv}<1>
            \begin{varblock}{example}{Creation}
                \begin{lstlisting}[style=MyBash]
                    |+Version 1.1.5 is the current production release
                    We are ready for "next release" ->+| |*NOW*| |+we decide v1.2+|
                    $ git switch -c release-1.2 develop
                    |+Switched to a new branch "release-1.2"+|
                    $ ./bump-version.sh 1.2  #Just modify some files
                    |+Files modified successfully, version bumped to 1.2.+|
                    $ git commit -a -m "Bump version number to 1.2"
                    |+[release-1.2 74d9424] Bumped version number to 1.2
                    1 files changed, 1 insertions(+), 1 deletions(-)+|
                \end{lstlisting}
            \end{varblock}
            \begin{varblock}{alerted}[0.8\textwidth]{}
                \alert{Adding new features here is \textbf{strictly prohibited}!}\\
                The idea is to just make sure everything works and is ready to be published!
            \end{varblock}
        \end{onlyenv}
        \begin{onlyenv}<2>
            \begin{varblock}{example}{Creation}
                \begin{lstlisting}[style=MyBash]
                    $ git switch -c release-1.2 develop
                    $ ./bump-version.sh 1.2 #Just modify some files
                    $ git commit -a -m "Bump version number to 1.2"
                \end{lstlisting}
            \end{varblock}
            \vspace{-2mm}
            \begin{varblock}{example}{Finalisation}
                \begin{lstlisting}[style=MyBash]
                    $ git switch main
                    |+Switched to branch 'main'+|
                    $ git merge --no-ff release-1.2
                    |+Merge made by recursive. (Summary of changes)+|
                    $ git tag -a v1.2
                    $ git switch develop
                    |+Switched to branch 'develop'+|
                    $ git merge --no-ff release-1.2  #Maybe conflicts!
                    |+Merge made by recursive. (Summary of changes)+|
                    $ git branch -d release-1.2
                    |+Deleted branch release-1.2 (was ff452fe).+|
                \end{lstlisting}
            \end{varblock}
        \end{onlyenv}
    \end{overlayarea}
\end{frame}
%---------------------------------------------------------------%

%---------------------------------------------------------------%
\subsubsection{Hotfix branches}
%---------------------------------------------------------------%
\begin{frame}[fragile,c]{The hotfix branch}
    \begin{tikzpicture}[remember picture, overlay, visible on=<2>]
        \node[inner sep=0pt] (yeah) at ([yshift=-8mm]current page.center) {
\includegraphics[height=40mm]{Meme-hotfix}};
        \node[below = 2mm of yeah, font=\scriptsize, text=PQ] {\ldots{}and, after all, never forget Murphy's law!};
        \node[above = 4mm of yeah, font=\Large, text=PT] {A hotfix is not something bad!};
        \draw[red, line width=1mm, shorten <= -1mm, shorten >= -1mm] (yeah.south west) -- (yeah.north east);
    \end{tikzpicture}
    \begin{onlyenv}<3>
        \vspace{-9mm}
        \begin{columns}[c]
            \begin{column}{0.55\textwidth}
                \begin{varblock}{}[\textwidth]{Properties:}
                    \begin{itemize}
                        \item May branch off from:\\
                        \quad{\small\main}
                        \item Must merge back into:\\
                        \quad{\small\develop,\;\main}
                        \item Branch naming convention:\\
                        \quad{\small\hotfix}
                    \end{itemize}
                \end{varblock}
                \vspace{-2mm}
                \begin{varblock}{alerted}{Main advantage}
                    \small
                    The essence is that work on the develop branch can continue, while a quick production fix is prepared
                \end{varblock}
            \end{column}
            \begin{column}{0.45\textwidth}
                \vspace{5mm}
                \centerline{\includegraphics[width=0.9\textwidth]{Hotfix}}
            \end{column}
        \end{columns}
    \end{onlyenv}
\end{frame}
%~~~~~~~~~~~~~~~~~~~~~~~~~~~~~~~~~~~~~~~~~~~~%
\begin{frame}[fragile,c]{The hotfix branch: How it works}
    \vspace{-6mm}
    \begin{overlayarea}{\textwidth}{0.8\textheight}
        \begin{onlyenv}<1>
            \begin{varblock}{example}{Creation}
                \begin{lstlisting}[style=MyBash]
                    $ git switch -c hotfix-1.2.1 main
                    |+Switched to a new branch "hotfix-1.2.1"+|
                    $ ./bump-version.sh 1.2.1  #Just modify some files
                    |+Files modified successfully, version bumped to 1.2.1.+|
                    $ git commit -a -m "Bump version number to 1.2.1"
                    |+[hotfix-1.2.1 41e61bb] Bump version number to 1.2.1
                    1 files changed, 1 insertions(+), 1 deletions(-)+|
                    # Some work to fix the situation
                    $ git commit -m "Fix severe production problem"
                    |+[hotfix-1.2.1 abbe5d6] Fix severe production problem
                    5 files changed, 32 insertions(+), 17 deletions(-)+|
                \end{lstlisting}
            \end{varblock}
        \end{onlyenv}
        \begin{onlyenv}<2>
            \begin{varblock}{example}{Creation}
                \begin{lstlisting}[style=MyBash]
                    $ git switch -c hotfix-1.2.1 main
                    $ ./bump-version.sh 1.2.1 #Just modify some files
                    $ git commit -a -m "Bump version number to 1.2.1"
                    # Some work to fix the situation
                    $ git commit -m "Fix severe production problem"
                \end{lstlisting}
            \end{varblock}
            \vspace{-3mm}
            \begin{varblock}{example}{Finalisation}
                \begin{lstlisting}[style=MyBash]
                    $ git switch main
                    |+Switched to branch 'main'+|
                    $ git merge --no-ff hotfix-1.2.1
                    |+Merge made by recursive. (Summary of changes)+|
                    $ git tag -a v1.2.1
                    $ git switch develop
                    |+Switched to branch 'develop'+|
                    $ git merge --no-ff hotfix-1.2.1
                    |+Merge made by recursive. (Summary of changes)+|
                    $ git branch -d hotfix-1.2.1
                    |+Deleted branch hotfix-1.2.1 (was abbe5d6).+|
                \end{lstlisting}
            \end{varblock}
        \end{onlyenv}
    \end{overlayarea}
\end{frame}
%---------------------------------------------------------------%

%---------------------------------------------------------------%
\subsubsection{Git-flow: an additional tool}
%---------------------------------------------------------------%
\begin{frame}[c]{Summary}
    \begin{tikzpicture}[remember picture, overlay]
        \node[anchor=west] (boss) at ([xshift=10mm]current page.west)
        {
\includegraphics[width=20mm, clip, trim=25mm 0 25mm 0]{Meme-like-a-boss}};
        \node[right = 8mm of boss]{\includegraphics[height=\textheight]{Git-model}};
    \end{tikzpicture}
\end{frame}
%~~~~~~~~~~~~~~~~~~~~~~~~~~~~~~~~~~~~~~~~~~~~%
\AlertFrame{%
    \begin{tabular}{c}
        
\includegraphics[width=3cm, clip, trim=0 50mm 0 20mm]{Meme-confused}\\
        But, wait, everything by hand?!\\[0.05\textheight]
        ~
    \end{tabular}
}
%~~~~~~~~~~~~~~~~~~~~~~~~~~~~~~~~~~~~~~~~~~~~%
\begin{frame}[fragile,c]{Git-flow: How to easily apply what we learnt}
    \vspace{-2mm}
    \begin{overlayarea}{\textwidth}{0.88\textheight}
        \begin{varblock}{example}[0.95\textwidth]{Pick a set of high-level repository operations}
            \PT{\href{https://github.com/nvie/gitflow}{\faGithub\;\textbf{nvie}}}
            \hspace{1ex}
            or better
            \hspace{1ex}
            \PP{\href{https://github.com/petervanderdoes/gitflow-avh}{\faGithub\;gitflow-avh}}
            \hspace{1ex}
            \PS{\href{https://github.com/CJ-Systems/gitflow-cjs}{\faGithub\;gitflow-cjs}}
        \end{varblock}
        \vspace{-1mm}
        \begin{onlyenv}<2>
            \centerline{\scriptsize \hspace{42mm}$\to\;$Standard raw git commands can be used as usual, too}
            \vspace{-6mm}
            \begin{varblock}{example}[0.9\textwidth]{The base syntax}
                \begin{lstlisting}[style=MyBash, aboveskip=2mm]
                    # To initialize a new repository, use:
                    git flow init |+[-d]+|
                    # To list/start/finish auxiliary branches, use:
                    git flow <type>
                    git flow <type> start <name> |+[<base>]+|
                    git flow <type> finish <name>
                    # To push/pull an auxiliary branch to the remote, use:
                    git flow <type> publish <name>
                    git flow <type> pull <remote> <name>
                \end{lstlisting}
                \vspace{1mm}
                \scriptsize
                \begin{tabular}{r@{\;}l}
                    \bash|<type>| & can be either \;\texttt{feature}, \;\texttt{release}\; or \;\texttt{hotfix}\\
                    \bash|<name>| & is the name of the auxiliary branch\\
                    \bash|<base>| & must be a commit on \;\texttt{develop}\; or \;\texttt{main}\\
                \end{tabular}
            \end{varblock}
        \end{onlyenv}
        \begin{onlyenv}<3>
            \begin{varblock}{alerted}[0.72\textwidth]{Take-home message}
                If you have a released software or plan to have one, Git-flow branching model can help you to bring your product to another level!
            \end{varblock}
            \begin{varblock}{}[\textwidth]{Think about it}
                Even if you use a (slightly) different branching model, it might well be that available git extensions still partially fit your use case
            \end{varblock}
        \end{onlyenv}
    \end{overlayarea}
    \FrameRemark<1-2>{Both AVH and CJS editions provide autocompletion for main shell types (checkout their README files for more information).}
\end{frame}
%~~~~~~~~~~~~~~~~~~~~~~~~~~~~~~~~~~~~~~~~~~~~%
\AlertFrame{I heard Git-flow is obsolete, is it?}
%~~~~~~~~~~~~~~~~~~~~~~~~~~~~~~~~~~~~~~~~~~~~%
\begin{frame}{No, it is not!}
    \PrepareURLsymbol[PB]
    \begin{onlyenv}<1>
        \begin{center}
            {\large\URL[PB]{https://trunkbaseddevelopment.com}{Trunk-based development}}\\[3mm]
            \includegraphics[width=\textwidth]{Git-trunk}\\[2mm]
            In the last years this alternative approach has become popular
        \end{center}
    \end{onlyenv}
    \begin{onlyenv}<2>
        \begin{varblock}{alert}[0.8\textwidth]{When does trunk-based development work best?}
            When you need to iterate quickly\\
            When you work mostly with senior developers
        \end{varblock}
        \begin{varblock}{example}[0.7\textwidth]{When does Git-flow work best?}
            When you have a lot of junior developers\\
            When you have an established product\\
            When you run an open-source project
        \end{varblock}
        \begin{varblock}{}[0.72\textwidth]{In academia}
            Development is rather slow and in small teams
        \end{varblock}
    \end{onlyenv}
    \FrameRemark<2>{If you are interested in a nice comparison of the two development techniques, refer to \URL*{https://www.toptal.com/software/trunk-based-development-git-flow}{this nice overview}.}
\end{frame}
%===============================================================%


%===============================================================%
\subsection{Summary and outlook}
%~~~~~~~~~~~~~~~~~~~~~~~~~~~~~~~~~~~~~~~~~~~~%
\begin{frame}{Most used Git commands}
    \PrepareURLsymbol[PT]
    \begin{tikzpicture}[remember picture, overlay]
        \node[anchor=north east, font=\small, draw=PT, thick, rounded corners]
            at ([shift={(-3mm,-2mm)}]current page.north east) {\URL*[PT]{https://git-scm.com/docs}{Git Documentation}};
        \begin{scope}[every node/.style={opacity=0.85}]
            \node (V) at (current page.center){
\includegraphics[width=0.85\textwidth]{V}};
            \node[anchor=south west] at (V.south west) {
\includegraphics[width=15mm]{V-logo}};
            \node[anchor=east, text=BGLIGHT] at ([xshift=-4mm]V.center) {
                \begin{tabular}{c}
                    There is\\
                    no certainty,\\
                    only opportunities!\\
                \end{tabular}
            };
        \end{scope}
    \end{tikzpicture}
\end{frame}
%===============================================================%

% Restore blue colours of Z02 theme for other sections
\colorlet{PB}{Marine}
\colorlet{PP}{SkyBlue}
\colorlet{PS}{BladeOfGrass}
\colorlet{PT}{CoralRed}
\colorlet{PQ}{Schizophrenia}
\colorlet{ALERT}{UnrealFoodPills}
\colorlet{BGDARK}{Aoi}
\colorlet{BGLIGHT}{BeachStorm!25}
\colorlet{TEXT}{Abyss}